% Generated draft: requires mathematical and source review.
\documentclass[11pt]{article}
\usepackage[margin=25mm]{geometry}
\usepackage{fontspec}
\setmainfont{Latin Modern Roman}
\usepackage{amsmath,amssymb,mathtools,unicode-math}
\setmathfont{Latin Modern Math}
\usepackage{enumitem,needspace,xurl,hyperref}
\setlength{\parindent}{0pt}
\setlength{\parskip}{5pt}
\setlength{\emergencystretch}{3em}
\newcommand{\N}{\mathbb{N}}
\newcommand{\Z}{\mathbb{Z}}
\newcommand{\Q}{\mathbb{Q}}
\newcommand{\R}{\mathbb{R}}
\newcommand{\C}{\mathbb{C}}
\newcommand{\F}{\mathbb{F}}
\newcommand{\A}{\mathbb{A}}
\newcommand{\PP}{\mathbb{P}}
\newcommand{\E}{\mathbb{E}}
\newcommand{\eps}{\varepsilon}
\newcommand{\dd}{\,\mathrm{d}}
\begin{document}
\section*{UnsolvedMath ID: 30005344}
\textbf{Source problem code:} catalogue entry
\subsection{Definitions and mathematical statement}
For a hypergraph \(H=(V,E)\), a matching is a set \(M\subseteq E\) of pairwise disjoint edges. If the edges carry colours given by a map \(c:E\to C\), then \(M\) is rainbow when its edges have pairwise distinct colours. Thus a rainbow matching of size \(t\) satisfies
\[
|M|=t,\qquad e\cap e'=\varnothing\quad\text{and}\quad c(e)\ne c(e')\qquad(e,e'\in M,\ e\ne e').
\]
A hypergraph is \(r\)-uniform if every edge has exactly \(r\) vertices. These definitions explain the terminology; the supplied record does not specify the full class of coloured hypergraphs intended in the question.

The requested problem is to determine how large \(N\) must be, in terms of \(t\) and \(r\), to guarantee a rainbow matching of size \(t\) in every admissible configuration. However, the supplied excerpt does not define \(N\), the admissible configurations, or their required colouring and matching hypotheses. Consequently, a fully quantified threshold statement cannot be recovered from this record alone.

The supplied dated curation reports that, for each fixed uniformity \(r\ge 3\), the functions denoted by \(f(r,t)\) and \(F(r,t)\) have order \(t^r\), with positive multiplicative constants depending on \(r\). It also reports that the conjecture attributed to Glebov--Sudakov--Szab\'o is false and describes that conjecture as linear. The record does not define these functions or state the conjecture's precise formulation. These are supplied contextual claims, not independently verified results in this draft. The curation labels the problem partially solved and says that the exact threshold remains undetermined.
\subsection{Short English statement}
Determine the threshold \(N\), as a function of \(t\) and \(r\), that forces \(t\) pairwise disjoint edges with distinct colours. The source excerpt omits what \(N\) counts and the hypotheses on the coloured hypergraph, so the intended threshold cannot yet be stated precisely. The supplied curation reports order \(t^r\) for fixed \(r\ge 3\), while describing the exact threshold as unresolved.
\subsection{Sources}
\begin{itemize}
\item Supplied local UnsolvedMath registry record, permanent ID \(30005344\), title: Exact Thresholds for Rainbow Matchings in Hypergraphs; exact source problem code: \texttt{catalogue entry}. The record places the problem in the 2023 Oberwolfach Reports collection and supplies no individual proposer.
\item Source report: Combinatorics (2023), Oberwolfach Reports, pp.~5--89, DOI 10.4171/owr/2023/1. \url{https://doi.org/10.4171/owr/2023/1}. The supplied background describes these reports as documenting mathematical workshops and their open-problem discussions. No precise page or passage locator for this question is supplied.
\item Supplied literature reference: Sharp Bounds for Rainbow Matchings. \url{https://arxiv.org/abs/2212.07580}.
\item Supplied literature reference: Sharp Bounds for Rainbow Matchings, JLMS. \url{https://doi.org/10.1112/jlms.70252}.
\item Attribution qualification: the excerpt names the conjecture of Glebov--Sudakov--Szab\'o and refers to an unidentified reference [10]. Its assertion beginning ``We showed'' does not identify the speakers in the supplied record.
\item Literature assessment checked at \texttt{2026-08-22}: the supplied curation labels the problem partially solved, reports order \(t^r\) for fixed \(r\ge 3\), reports that the linear conjecture is false, and describes the exact threshold as undetermined. This is supplied dated curation derived from the OWR material, not independently verified literature assessment; the record requests independent verification of time-sensitive status and cited claims.
\item The registry labels the statement \texttt{corrected\_verified} and says it was checked against the official report and preserves its hypotheses, notation, and target. That verification claim is supplied metadata; this draft has not independently checked the report, and the supplied excerpt still lacks the hypotheses needed for a precise threshold formulation.
\end{itemize}
\end{document}
